\documentclass[10pt,a4paper]{amsart}
\usepackage[margin=19mm]{geometry}
\usepackage{amsmath,amssymb,mathtools}
\usepackage[scr=rsfs,cal=euler]{mathalfa}
\usepackage{xfrac}
\usepackage[unicode,hidelinks,bookmarks=false]{hyperref}
\newcommand{\ie}{i.e.\ }
\newcommand{\cf}{cf.\ }
\newcommand{\resp}{respectively\ }
\newcommand{\Subsection}{\subsection}
\providecommand{\nobreakdash}{\nobreak\hbox{-}}
\setlength{\emergencystretch}{2em}
\multlinegap0pt
\usepackage{xcolor}
\usepackage{mathtools}
\usepackage{algorithm}\usepackage{algpseudocode}
\usepackage{enumerate}
\usepackage{bm}
\newtheorem{lemma}{Lemma}
\usepackage{cleveref}
\crefname{lemma}{Lemma}{Lemmas}
\newcommand{\R}{\mathbb{R}}
\DeclareMathOperator*{\argmin}{arg\,min}
\def\evu{{u}}
\newcommand{\nn}{\nonumber} % Nonumber in equation
\def\vv{{\bm{v}}}
\title{Revisiting Incremental Stochastic Majorization-Minimization Algorithms with Applications to Mixture of Experts}
\author{TrungKhang Tran, TrungTin Nguyen, Gersende Fort, Tung Doan, Hien Duy Nguyen, Binh T. Nguyen, Florence Forbes, Christopher Drovandi}
\date{Source: arXiv:2601.19811v1 (CC BY 4.0)}
\begin{document}
\maketitle

\noindent\emph{Source and licence.} Revisiting Incremental Stochastic Majorization-Minimization Algorithms with Applications to Mixture of Experts. Version arXiv:2601.19811v1 (CC BY 4.0), distributed by arXiv under the Creative Commons Attribution 4.0 International licence, http://creativecommons.org/licenses/by/4.0/, as recorded in the arXiv licence metadata for this exact version and shown on the abstract page https://arxiv.org/abs/2601.19811v1. Full licence notice: \texttt{licenses/arxiv-2601.19811-CC-BY-4.0.txt}.\par\smallskip

\noindent\textit{Editorial scope.} Counted unit: the article's lemma lemm\_notes\_Upsilon\_Sigma (source0.tex lines 3060--3067). The article's complete original proof of it is delivered with it (source0.tex lines 3068--3079, one \texttt{proof} environment of 12 lines). No mathematics is written, repaired or re-derived: the statement and the proof are the article's own lines, in the article's own notation, and no retained line is substituted or re-flowed.  No further source-local context is needed: every cross-reference of every retained line resolves inside the two delivered spans.  Nothing was omitted from any retained span, so the package carries no omission record.  No retained line contains a citation, so the wrapper carries no bibliography and no bibliography entry.  No retained line contains a figure environment, an image-inclusion command or drawing code, so no image is delivered and the package declares no asset dependency.  Editorial formatting only, in addition: an AMS \texttt{amsart} preamble that restates the article's own declarations and macros used by the retained lines, namely the environments \texttt{lemma} on separate counters, together with the standard packages those lines need.  The retained environments print in this document as Lemma 1 in order of appearance, whereas the article prints the same items with the numbering of its own full document.  The complete native source of the article as distributed in the arXiv e-print for this exact version is bundled unchanged as \texttt{sources/193488/source0.tex}.
\begin{lemma} \label{lemm_notes_Upsilon_Sigma}
    Consider the function 
    \begin{align}
        f(\evu,\vv)=\frac{g(\vv)}{\evu}+c\log(\evu). \nn
    \end{align}
    Where $\vv$ is a vector, $\evu$ is a positive real number, $c$ is a positive constant and $g(\vv): \R^{n}\rightarrow \R^{+}$ be a function that has exactly one global minimum. 
    Denote $\vv_0=\displaystyle\argmin_{\vv}g(\vv)$ be the global minimum of $g(\vv)$ and $\evu_0=\displaystyle\argmin_{\evu}\left(\frac{g(\vv_0)}{\evu}+\log(\evu)\right)$. We have $f(\evu_0,\vv_0)$ is the global minimum of $f$ over all $\evu,~\vv$.
\end{lemma}

\begin{proof}[Proof of \cref{lemm_notes_Upsilon_Sigma}]
    For all $\vv$, we have:
    \begin{align*}
        &\nabla_u f(\evu,\vv)=\frac{1}{\evu^2}(-g(\vv)+cu),~\nabla^2_u f(\evu,\vv)=\frac{1}{\evu^3}(2g(\vv)-cu).
    \end{align*}
    Hence $\evu=\frac{g(\vv)}{c}$ is the local minimum of $f(\evu,\vv$) for all $\vv$. But since $\nabla_{\evu}f(\evu,\vv)$ has exactly one solution for $\evu$ and both $\displaystyle\lim_{\evu\rightarrow+\infty}f(\evu,\vv), \displaystyle\lim_{\evu\rightarrow+0}f(\evu,\vv)$ is $+\infty$. We can conclude that $\evu=\frac{g(\vv)}{c}$ is the global minimum of $f(\evu,\vv$) for all given $\vv$ .\\
    Assume that $(\evu_1,\vv_1)$ is the minimum of $f(\evu,\vv$), we have
    \begin{align*}
        f(\evu_1,\vv_1)\geq f\left(\frac{g(\vv_1)}{c},\vv_1\right)=c\log\left(\frac{g(\vv_1)}{c}\right) \geq c\log\left(\frac{g(\vv_0)}{c}\right)= f\left(\frac{g(\vv_0)}{c},\vv_0\right).
    \end{align*}
    \par Which lead to $g(\vv_1)=g(\vv_0)$. But since  $\vv_0$ is the only global minimum of $g(\vv)$, we must have $\vv_1=\vv_0$. Thus, leading to $\evu_1$ also is equal to $\evu_0$.
\end{proof}

\end{document}
